\section{Objectives of constrained control}

In the previous chapter we formulated the constrained finite time optimal control problem that lies
at the core of nominal MPC. However, solving a finite-horizon problem at one instant is only part
of the story. In closed loop, one wants the controller to keep generating admissible inputs over time,
so that the state and input constraints are never violated. This naturally leads to the question of
which initial states are \emph{safe} in the sense that constraint satisfaction can be maintained forever.\\

Consider a general discrete-time system
\[
x(k+1)=g(x(k),u(k)),
\]
with state and input constrained to lie in sets
\[
x(k)\in \cX,\qquad u(k)\in \cU.
\]
A control law \(u(k)=\kappa(x(k))\) should ideally achieve several goals at once. First, it should
respect the constraints for all time. Second, it should stabilize the system, typically to the origin
or to a desired equilibrium. Third, among all admissible closed-loop behaviors, it should optimize
a suitable performance criterion. Finally, one would like the controller to work from as large a set
of initial conditions as possible.\\

In this chapter we focus on the first of these goals: persistent constraint satisfaction. This is the
role of invariant and control invariant sets. Roughly speaking, an invariant set describes a region in
which the system can remain forever once it enters it, while a control invariant set describes a region
from which one can always choose an admissible input that keeps the system inside that region.
These notions are fundamental in MPC because they formalize the idea of long-term feasibility.

A simple way to understand why such sets matter is to compare linear feedback with predictive
control in the presence of hard constraints. Even if an unconstrained linear controller stabilizes the
system, its region of safe operation may be very small once actuator limits are imposed. Outside
that region, either the input saturates immediately or the state evolves in such a way that future
constraint violations become unavoidable. Nonlinear constrained controllers, and in particular MPC,
are attractive precisely because they can enlarge the set of initial states from which the constraints
can be respected indefinitely.

\section{Invariance}

\subsection{Positively invariant sets}

We begin with the autonomous case. Consider the system
\[
x(k+1)=g(x(k)),
\]
or, more generally, a closed-loop system obtained after fixing a feedback law \(\kappa\),
\[
x(k+1)=g(x(k),\kappa(x(k))).
\]
The central question is the following: for which initial states does the trajectory remain inside the
constraint set \(\cX\) for all future times?

\begin{theorembox}[Positively invariant sets]
A set \(\mathcal O\subseteq \R^n\) is called a \emph{positively invariant set} for the autonomous system
\(x(k+1)=g(x(k))\) if
\[
x(k)\in \mathcal O \quad \Longrightarrow \quad x(k+1)\in \mathcal O
\qquad \forall k\in \{0,1,2,\dots\}.
\] 
\end{theorembox}

Thus, once the state starts inside \(\mathcal O\), all future states remain inside \(\mathcal O\) as well.
If in addition \(\mathcal O\subseteq \cX\), then \(\mathcal O\) provides a set of initial conditions from
which the trajectory never violates the state constraints.\\
Among all positively invariant sets contained in the constraint set \(\cX\), one is especially important.

\begin{theorembox}[Maximal positively invariant set]
The set \(\mathcal O_\infty\subseteq \cX\) is called the \emph{maximal positively invariant set} with
respect to \(\cX\) if
\[
\mathcal O_\infty \text{ is positively invariant}
\]
and
\[
\mathcal O \subseteq \mathcal O_\infty
\qquad
\text{for every positively invariant set } \mathcal O \subseteq \cX.
\]
\end{theorembox}

This set has a very natural interpretation: it is the set of all states for which the autonomous
system can remain feasible forever once it is there. In that sense, \(\mathcal O_\infty\) is the largest
region of perpetual constraint satisfaction for the given dynamics.

\subsection{Pre-sets}

To study invariance, it is useful to introduce a one-step backward reachability notion.

\begin{theorembox}[Pre-set]
Given a set \(S\) and the dynamic system
\[
x(k+1)=g(x(k)),
\]
the \emph{pre-set} of \(S\) is the set of states that evolve into the target set \(S\) in one time step:
\[
\operatorname{pre}(S):=\{x \mid g(x)\in S\}.
\]
\end{theorembox}

So, if \(S\) is the set of states we would like to be in at the next time instant, then
\(\operatorname{pre}(S)\) is the set of states from which this happens in one step. In other words,
the pre-set tells us which current states are compatible with reaching the target set immediately
under the given dynamics.

This notion gives a very compact way to express invariance. Indeed, a set \(\mathcal O\) is positively
invariant precisely when every state in \(\mathcal O\) is mapped back into \(\mathcal O\) after one step.
That means exactly that every state in \(\mathcal O\) must belong to the pre-set of \(\mathcal O\).

\begin{theorembox}[Geometric condition for invariance]
A set \(\mathcal O\) is a positively invariant set if and only if
\[
\mathcal O \subseteq \operatorname{pre}(\mathcal O).
\]
Equivalently,
\[
\operatorname{pre}(\mathcal O)\cap \mathcal O = \mathcal O.
\]
\end{theorembox}
The proof is immediate and is conveniently written through the contrapositive in both directions.
If \(\mathcal O \not\subseteq \operatorname{pre}(\mathcal O)\), then there exists some
\(\bar x\in \mathcal O\) such that \(\bar x\notin \operatorname{pre}(\mathcal O)\). By definition of
the pre-set, this means \(g(\bar x)\notin \mathcal O\), so \(\mathcal O\) is not positively invariant.
Conversely, if \(\mathcal O\) is not positively invariant, then there exists some \(\bar x\in \mathcal O\)
such that \(g(\bar x)\notin \mathcal O\). Again by definition of the pre-set, this implies
\(\bar x\notin \operatorname{pre}(\mathcal O)\), and therefore
\(\mathcal O\not\subseteq \operatorname{pre}(\mathcal O)\).\\

For general nonlinear systems, computing pre-sets can be very difficult. Even if the target set \(S\)
has a simple description, the set of all states that reach it in one time step may have a complicated
shape. Only in special cases, most notably for linear systems with suitably described constraint
sets, does the computation become tractable.

\subsection{Linear autonomous systems}

For linear autonomous dynamics
\[
x(k+1)=Ax(k),
\]
the pre-set becomes much easier to describe. If the target set is a polyhedron
\[
S=\{x\in \R^n \mid Fx\leq f\},
\]
then
\[
\operatorname{pre}(S)
=
\{x\in \R^n \mid Ax\in S\}
=
\{x\in \R^n \mid FAx\leq f\}.
\]
So for linear systems the pre-set of a polyhedron is again a polyhedron. This closure property is
the key reason why polyhedral invariant-set computations are possible.

\subsection{Conceptual computation of the maximal invariant set}

The maximal positively invariant set can be characterized by a simple iterative construction.
Starting from the full constraint set,
\[
\Omega_0 := \cX,
\]
one recursively removes all states that cannot remain in the current set after one time step. This
leads to the iteration
\[
\Omega_{i+1}:=\operatorname{pre}(\Omega_i)\cap \Omega_i.
\]
Since \(\Omega_{i+1}\) is obtained by intersecting \(\Omega_i\) with the set of states that return to
\(\Omega_i\) in one step, the sequence is nested:
\[
\Omega_{i+1}\subseteq \Omega_i, \qquad \forall i\in\mathbb N.
\]
Thus, at each step, one removes the states that cannot be kept inside the current candidate set.\\
The interpretation is very natural. At the first iteration one discards all states in \(\cX\) that leave
\(\cX\) after one transition. At the second iteration one discards all states that enter the already
discarded region after one transition, and so on. If the procedure reaches a fixed point, that fixed
point is precisely the maximal positively invariant set.

\begin{theorembox}[Conceptual algorithm to compute the maximal invariant set]
\textbf{Input:} \(g,\cX\)\\
\textbf{Output:} \(\mathcal O_\infty\)

\[
\begin{array}{l}
\Omega_0 \leftarrow \cX \\[0.3em]
\textbf{loop} \\[0.2em]
\qquad \Omega_{i+1} \leftarrow \operatorname{pre}(\Omega_i)\cap \Omega_i \\[0.2em]
\qquad \textbf{if } \Omega_{i+1}=\Omega_i \textbf{ then} \\[0.2em]
\qquad\qquad \textbf{return } \mathcal O_\infty=\Omega_i \\[0.2em]
\qquad \textbf{end if} \\[0.2em]
\textbf{end loop}
\end{array}
\]
\end{theorembox}

So, if for some finite index \(i\) the iteration satisfies
\[
\Omega_{i+1}=\Omega_i,
\]
then the algorithm has converged and the limit set is
\[
\mathcal O_\infty=\Omega_i.
\]
This set is the largest positively invariant subset of \(\cX\).\\
In practice, each iteration requires several set operations: one must represent the current set
\(\Omega_i\), compute its pre-set, intersect sets, and check whether the procedure has converged.
These operations are manageable in some structured settings, especially for linear systems with
polyhedral sets, but they can become much more difficult in general nonlinear problems.

\section{Control invariance}

\subsection{Definition and meaning}

In the autonomous setting, the evolution is fixed once the current state is known. In control, by
contrast, we can influence the next state through the input. This leads to a more flexible notion.\\
Consider the controlled system
\[
x(k+1)=g(x(k),u(k)),
\]
subject to the constraints
\[
x(k)\in \cX,\qquad u(k)\in \cU.
\]
\begin{theorembox}[Control invariant sets]
A set \(\mathcal C\subseteq \cX\) is called a \emph{control invariant set} if
\[
x(k)\in \mathcal C
\quad \Longrightarrow \quad
\exists\,u(k)\in \cU \text{ such that } g(x(k),u(k))\in \mathcal C
\]
for all \(k\in \mathbb N\).
\end{theorembox}

This definition differs from positive invariance in one essential way: it does not require that the
next state remains in the set under a fixed autonomous evolution, but only that there \emph{exists}
an admissible control action that keeps the successor state in the set. Therefore control invariance
captures the states from which constraint satisfaction can be maintained indefinitely by choosing
the input appropriately.\\
The largest such set plays the same role as before.

\begin{theorembox}[Maximal control invariant set]
The set \(\mathcal C_\infty\subseteq \cX\) is called the \emph{maximal control invariant set} for the
system \(x(k+1)=g(x(k),u(k))\) under the constraints \((x,u)\in \cX\times \cU\) if it is control
invariant and contains every control invariant set contained in \(\cX\).
\end{theorembox}

Thus \(\mathcal C_\infty\) is the best any controller can do in terms of perpetual constraint
satisfaction. If \(x(0)\in \mathcal C_\infty\), then there exists at least one admissible control strategy
that keeps the system feasible forever. If \(x(0)\notin \mathcal C_\infty\), then no controller can
guarantee this.

\subsection{Controlled pre-sets}

The pre-set concept extends naturally to systems with inputs. Given a target set \(S\subseteq \R^n\),
we define
\[
\operatorname{pre}(S):=
\{x\in \R^n \mid \exists\,u\in \cU \text{ such that } g(x,u)\in S\}.
\]
So the pre-set now contains all states from which it is \emph{possible} to drive the system into \(S\)
in one step using an admissible input.\\
With this definition, the geometric condition for control invariance has exactly the same form as
before.

\begin{theorembox}[Geometric condition for control invariance]
A set \(\mathcal C\) is control invariant if and only if
\[
\mathcal C \subseteq \operatorname{pre}(\mathcal C).
\]
\end{theorembox}

Accordingly, the same conceptual fixed-point iteration can be used:
\[
\Omega_0:=\cX,
\qquad
\Omega_{i+1}:=\operatorname{pre}(\Omega_i)\cap \Omega_i.
\]
If the iteration converges in finite time, the fixed point is the maximal control invariant set
\(\mathcal C_\infty\).\\

Conceptually this is identical to the autonomous case, but computationally it is much harder,
because the controlled pre-set involves existential quantification over the input. In other words,
one must determine not only where the state goes, but whether there exists an admissible input
that drives it into the target set.

\subsection{Linear systems and projection}

For the linear system
\[
x(k+1)=Ax(k)+Bu(k),
\]
with input constraints
\[
\cU=\{u\in \R^m \mid Gu\leq g\},
\]
and target set
\[
S=\{x\in \R^n \mid Fx\leq f\},
\]
the controlled pre-set is
\[
\operatorname{pre}(S)
=
\{x\in \R^n \mid \exists\,u\in \cU \text{ such that } Ax+Bu\in S\}.
\]
Equivalently,
\[
\operatorname{pre}(S)
=
\left\{
x\in \R^n \,\middle|\,
\exists\,u\in \R^m:
\begin{bmatrix}
FA & FB\\
0 & G
\end{bmatrix}
\begin{bmatrix}
x\\u
\end{bmatrix}
\leq
\begin{bmatrix}
f\\g
\end{bmatrix}
\right\}.
\]
This description shows that the controlled pre-set is obtained by projecting a higher-dimensional
polyhedron in the \((x,u)\)-space onto the \(x\)-space. That projection step is the main source of
computational difficulty.\\
Still, the conceptual meaning is very clear. In the autonomous case, each state has only one
possible successor. In the controlled case, each state is associated with an entire set of possible
successors, one for each admissible input. This makes the controlled pre-set much larger, but also
much more difficult to compute.

\subsection{From a control invariant set to a control law}

A control invariant set is not just a geometric object: it can be used directly to synthesize a control
law that guarantees constraint satisfaction. Suppose \(\mathcal C\) is control invariant. Then for every
\(x\in \mathcal C\) there exists at least one input \(u\in \cU\) such that the next state remains in
\(\mathcal C\). One can therefore define a feedback law by selecting, at each state, one admissible
input with this property. A generic construction is
\[
\kappa(x)
:=
\argmin_{u}
\left\{
f(x,u)\ \middle|\ g(x,u)\in \mathcal C,\ u\in \cU
\right\},
\]
where \(f\) can be any objective function, even \(f(x,u)=0\) if one only wants feasibility.\\
If \(\kappa(x)\in \cU\) for all \(x\in \mathcal C\) and
\[
g(x,\kappa(x))\in \mathcal C \qquad \forall x\in \mathcal C,
\]
then the closed-loop system remains in \(\mathcal C\) forever, and therefore the constraints are
satisfied for all time.\\
This construction guarantees feasibility, but not necessarily convergence. Stability requires
additional ingredients. This is precisely where MPC enters: it augments feasibility with an optimal
control objective and suitable terminal conditions so as to recover both constraint satisfaction and
closed-loop stability.

\section{Relation to MPC}

Invariant and control invariant sets are fundamental in MPC, but one rarely computes the maximal
control invariant set explicitly in realistic problems. The reason is not conceptual but computational.
Even for constrained linear systems, exact computation of maximal control invariant sets can become
very complex. For nonlinear systems it is usually far harder.\\

This difficulty explains one of the deep ideas behind MPC. Instead of trying to compute the largest
control invariant region once and for all, MPC solves a tractable finite-horizon optimization problem
at each sampling instant. In doing so, it implicitly describes a control invariant behavior in a form
that is much easier to represent and compute online. Terminal ingredients such as invariant terminal
sets are then used to connect the finite-horizon problem with the infinite-horizon feasibility and
stability properties one would ideally like to have.

\begin{theorembox}[Control invariance and MPC]
A maximal control invariant set tells us the largest region from which some admissible controller can
keep the constraints satisfied forever. MPC can be viewed as a tractable way of realizing this idea
indirectly, without having to compute that maximal set explicitly.
\end{theorembox}

\section{Practical computation of invariant sets}

\subsection{Polyhedra and polytopes}

For computational purposes, invariant sets are often represented as polyhedra. A polyhedron is the
intersection of finitely many halfspaces and can be written as
\[
P=\{x\in \R^n \mid Ax\leq b\}.
\]
If \(P\) is bounded, it is called a \emph{polytope}. In linear MPC this is the natural set class,
because state, input, and terminal constraints are often given by linear inequalities.\\
A complementary viewpoint is given by the convex hull representation. If \(v_1,\dots,v_s\in \R^n\),
their convex hull is
\[
\operatorname{co}\{v_1,\dots,v_s\}
=
\left\{
x=\sum_{i=1}^s \lambda_i v_i
\ \middle|\
\lambda_i\geq 0,\ \sum_{i=1}^s \lambda_i=1
\right\}.
\]
The Minkowski--Weyl theorem states that a set is a polytope if and only if it can be represented
both as a bounded intersection of finitely many halfspaces and as the convex hull of finitely many
points. In computations, one switches between these two descriptions depending on which set
operation must be performed.

\subsection{Common constraints in MPC}

Many constraints that appear in MPC are naturally polytopic. Simple box constraints on inputs and
outputs,
\[
u_{\mathrm{low}}\leq u(k)\leq u_{\mathrm{high}},
\qquad
y_{\mathrm{low}}\leq y(k)\leq y_{\mathrm{high}},
\]
can be written as linear inequalities in the stacked vector \(\begin{bmatrix}x(k)\\u(k)\end{bmatrix}\).
Likewise, magnitude constraints such as
\[
\norm{Cx(k)}_\infty \leq \alpha
\]
and rate constraints such as
\[
\norm{x(k+1)-x(k)}_\infty \leq \alpha
\]
are also polyhedral. This is why polytopes are such a natural computational language in linear
MPC.

\subsection{Intersection and subset tests}

The iterative computation of invariant sets relies repeatedly on two basic operations: intersection
and equality testing.

If
\[
S=\{x\in \R^n \mid Cx\leq c\},
\qquad
T=\{x\in \R^n \mid Dx\leq d\},
\]
then their intersection is obtained simply by stacking inequalities:
\[
S\cap T
=
\left\{
x\in \R^n \,\middle|\,
\begin{bmatrix}
C\\D
\end{bmatrix}x
\leq
\begin{bmatrix}
c\\d
\end{bmatrix}
\right\}.
\]
So intersection is easy in inequality form.

Testing whether one polytope is contained in another is more subtle. Suppose
\[
P=\{x\in \R^n \mid Cx\leq c\},
\qquad
Q=\{x\in \R^n \mid Dx\leq d\}.
\]
To check whether \(P\subseteq Q\), it is enough to verify each inequality defining \(Q\).
Let \(D_i\) denote the \(i\)-th row of \(D\). Then \(P\subseteq Q\) holds if and only if
\[
\max_{x}\ D_i x
\quad \text{subject to } Cx\leq c
\]
is at most \(d_i\) for every row \(i\). This maximum is the support function of \(P\) in the
direction \(D_i\):
\[
h_P(D_i):=\max_x\ D_i x \quad \text{subject to } Cx\leq c.
\]
Therefore
\[
P\subseteq Q
\qquad \Longleftrightarrow \qquad
h_P(D_i)\leq d_i \quad \text{for all } i.
\]

This subset test is important because equality of two polytopes can be checked through two
containment relations:
\[
P=Q
\qquad \Longleftrightarrow \qquad
P\subseteq Q \text{ and } Q\subseteq P.
\]

\subsection{Polytopic invariant sets}

For linear autonomous systems with polytopic constraints, the conceptual invariant-set algorithm
becomes implementable. One represents each iterate \(\Omega_i\) as a polytope, computes the
pre-set through linear inequality manipulation, intersects polytopes by stacking inequalities, and
tests convergence through bidirectional subset checks.

In principle, this allows one to obtain the maximal polyhedral invariant set. In practice, however,
the number of inequalities can grow quickly with the state dimension and with the number of
iterations. This is one of the main limitations of exact polyhedral computations in higher
dimensions.

\subsection{Ellipsoids and Lyapunov-based invariant sets}

An alternative is to work with ellipsoids. Given a symmetric positive definite matrix \(P\in \R^{n\times n}\)
and a center \(x_c\in \R^n\), the ellipsoid
\[
\mathcal E:=\{x\in \R^n \mid (x-x_c)^\top P (x-x_c)\leq 1\}
\]
has a fixed-complexity description: no matter how large the dimension is, containment can be
checked through quadratic expressions rather than through arbitrarily many inequalities.

This connects naturally with Lyapunov theory. Suppose \(V:\R^n\to \R\) is a Lyapunov function
for the autonomous system
\[
x(k+1)=g(x(k)).
\]
Then every sublevel set
\[
Y_\alpha:=\{x\in \R^n \mid V(x)\leq \alpha\}
\]
is invariant. Indeed, if the state starts in \(Y_\alpha\), then the Lyapunov function decreases along
the trajectory, so the state can never leave that sublevel set.

For linear systems
\[
x(k+1)=Ax(k),
\]
a standard quadratic Lyapunov function is
\[
V(x)=x^\top P x,
\]
where \(P\succ 0\) satisfies
\[
A^\top P A - P \prec 0.
\]
Then the invariant sublevel sets are ellipsoids of the form
\[
Y_\alpha=\{x\in \R^n \mid x^\top P x\leq \alpha\}.
\]
Such sets are typically not maximal, unlike the exact polyhedral invariant set, but they are much
easier to compute and scale better to larger dimensions.

\subsection{Largest ellipsoid contained in a polytope}

Suppose the state constraints are given by the polytope
\[
\cX=\{x\in \R^n \mid Fx\leq f\},
\]
and suppose \(V(x)=x^\top P x\) is a Lyapunov function. We want the largest ellipsoidal invariant
set of the form
\[
Y_\alpha=\{x\in \R^n \mid x^\top P x\leq \alpha\}
\]
that is contained in \(\cX\).

This means choosing the largest \(\alpha\geq 0\) such that
\[
Y_\alpha \subseteq \cX.
\]
Using the support function, this is equivalent to requiring
\[
h_{Y_\alpha}(F_i)\leq f_i
\qquad \text{for all rows } F_i \text{ of } F.
\]
For an ellipsoid one can compute this support function explicitly:
\[
h_{Y_\alpha}(F_i)=\norm{P^{-1/2}F_i^\top}\sqrt{\alpha}.
\]
Hence the inclusion condition becomes
\[
\norm{P^{-1/2}F_i^\top}\sqrt{\alpha}\leq f_i
\qquad \forall i,
\]
or, equivalently,
\[
\alpha \leq \frac{f_i^2}{F_i P^{-1}F_i^\top}
\qquad \forall i.
\]
Therefore the largest admissible value is
\[
\alpha^\star
=
\min_{i}
\frac{f_i^2}{F_i P^{-1}F_i^\top}.
\]
This gives the largest Lyapunov ellipsoid of the chosen shape matrix \(P\) that is still fully
contained in the constraints.

The important message is that ellipsoidal invariant sets trade exactness for tractability. They are
usually smaller than the maximal polyhedral invariant set, but they are easy to evaluate and can be
computed efficiently even in higher dimensions.

\section{Summary}

Invariant and control invariant sets provide the right geometric language for infinite-horizon
constraint satisfaction. A positively invariant set describes a region in which an autonomous or
already closed-loop system remains forever once it enters it. A control invariant set describes a
region in which one can always choose an admissible input that keeps the state inside the region.

The maximal positively invariant set \(\mathcal O_\infty\) and the maximal control invariant set
\(\mathcal C_\infty\) are the largest regions with these properties. They can be characterized through
pre-sets and, at least conceptually, computed by the same fixed-point iteration
\[
\Omega_{i+1}=\operatorname{pre}(\Omega_i)\cap \Omega_i.
\]
For autonomous systems this iteration removes states that cannot remain feasible after one step.
For controlled systems it removes states from which no admissible input can keep the trajectory
feasible.

In linear settings with polyhedral constraints, invariant-set computations can be carried out using
polytopes. Intersections are easy in inequality form, while subset tests can be performed through
support functions. Exact polyhedral invariant sets are very informative, but their complexity can
grow quickly with dimension.

Ellipsoidal invariant sets offer a more tractable alternative. When obtained from quadratic
Lyapunov functions, they are invariant by construction and their size inside a polytopic constraint
set can be computed explicitly. They are generally conservative compared with maximal polyhedral
sets, but their complexity remains fixed and manageable.

From the viewpoint of MPC, the key lesson is that invariant sets are not just auxiliary geometric
objects. They are central tools for reasoning about long-term feasibility, and they form the bridge
between finite-horizon optimization and infinite-horizon constraint satisfaction. In the next chapter,
this bridge will be used to study recursive feasibility and stability of MPC.